We develop a framework for assessing the economic consequences of national resource
plans. The framework links CLEWS, a planning model that identifies the least-cost way to
meet a country's energy, land and water needs, with OG-Core, an overlapping-generations
model of the economy in which households differ by age and lifetime income. Eight channels
carry electricity costs, infrastructure investment, changes in generation technology, air
pollution and policy instruments into the economy, and return the economy's demand and
cost of capital to resource planning. We use a Philippine scenario to show the kinds of
information the framework provides: the size of each effect, how effects unfold over time,
and how a change in one part of the economy spreads to others through prices, wages and
returns. More broadly, the framework allows a resource plan to be assessed for its effects
on production, public finances and the distribution of costs and benefits across
generations, not only for its cost and feasibility.
